% !TeX root = ../../Book4.tex
% 纲要标题：### 17.2 三角范畴公理
% 纲要标签：b4:sec:1602
% 纲要位置：计划纲要.md，第 3661 行。

\section{三角范畴公理}
\label{b4:sec:1602}

映射锥的旋转和复合规律可以独立表述为三角范畴公理。核对八面体时要保留各箭头的符号，才能将抽象相容性落实到复形公式。


% 纲要内容（保留纲要核验项目）：
%
% * 好三角与旋转
% * 态射公理与八面体公理
% * 三角范畴不是一般交换范畴
% * 同伦范畴的三角结构通过复形层面的显式锥构造验证；八面体公理的矩阵映射与符号单独核对，不把其成立仅列作形式约定
%

\subsection{好三角与旋转}
\label{b4:subsec:1602:01}

对可复合的 \(X\xrightarrow{u}Y\xrightarrow{v}Z\)，锥 \(C_u\) 记录第一步的差别，\(C_v\) 记录第二步的差别，而 \(C_{vu}\) 记录整个复合的差别。八面体公理把三者联系为一条好三角
\[
 C_u\longrightarrow C_{vu}\longrightarrow C_v\longrightarrow C_u[1].
\]
三角范畴将这种联系与锥的补全、旋转及态射补全一并作为公理，因此不必要求对象本身仍是复形。

\begin{definition}\label{b4:def:triangulated-category}
设 \(\mathcal T\) 为加性范畴，给定自等价 \([1]\) 及其逆，并固定它们的迭代。再指定一类称为好三角的图式 \(X\xrightarrow{u}Y\xrightarrow{v}Z\xrightarrow{w}X[1]\)。若这些数据满足下面四条公理，则称 \(\mathcal T\) 连同这些数据为\term[sanjiao-fanchou]{三角范畴}[triangulated category]。
\begin{enumerate}
\item[(TR1)] 好三角对同构封闭；\(X\xrightarrow{1}X\to0\to X[1]\) 是好三角；每个态射 \(u\) 都能补成好三角。
\item[(TR2)] 一个三角为好三角，当且仅当其旋转
\(Y\xrightarrow{v}Z\xrightarrow{w}X[1]\xrightarrow{-u[1]}Y[1]\) 为好三角。
\item[(TR3)] 两个好三角的前两项之间若已有交换方块，则可以补上第三项的态射，成为三角态射。
\item[(TR4)] 对可复合的 \(X\xrightarrow{u}Y\xrightarrow{v}Z\)，取分别补全 \(u,v,vu\) 的好三角，其后两条箭头记为 \(i_u,p_u\)、\(i_v,p_v\)、\(i_{vu},p_{vu}\)，第三对象记为 \(C_u,C_v,C_{vu}\)。存在 \(\alpha:C_u\to C_{vu}\)、\(\beta:C_{vu}\to C_v\)，满足
\[
\alpha i_u=i_{vu}v,\quad p_{vu}\alpha=p_u,\quad
\beta i_{vu}=i_v,\quad p_v\beta=u[1]p_{vu},
\]
并使 \(C_u\xrightarrow{\alpha}C_{vu}\xrightarrow{\beta}C_v\xrightarrow{i_u[1]p_v}C_u[1]\) 为好三角。
\end{enumerate}
\end{definition}

旋转中的负号不能任意删去。它与平移复形的负微分一起，使同一条长正合列在换起点后仍由三角产生。连续旋转三次所得是三个箭头都带负号的平移三角；在三个对象上分别取 \(1,-1,1\) 作同构，可将它改写为箭头 \(u[1],v[1],-w[1]\) 的好三角。这里最后一条箭头仍带负号，不能把三个符号一起删去。

\ProseParagraphBreak
以下先从前三条公理推导 Hom 正合性及其应用，最后用锥的具体计算验证 \(K(\mathcal A)\) 满足全部公理。这样，在验证八面体时调用的正合性只依赖已经核对的前三条公理。

\subsection{态射公理与八面体公理}
\label{b4:subsec:1602:02}

(TR3) 只要求补全存在，没有要求唯一。(TR4) 则规定这种补全对于复合态射还应满足什么相容关系。把三个三角写成
\[
 X\xrightarrow{u}Y\xrightarrow{i_u}C_u\xrightarrow{p_u}X[1],\quad
 Y\xrightarrow{v}Z\xrightarrow{i_v}C_v\xrightarrow{p_v}Y[1],
\]
\[
 X\xrightarrow{vu}Z\xrightarrow{i_{vu}}C_{vu}\xrightarrow{p_{vu}}X[1].
\]
所要求的是 \(\alpha:C_u\to C_{vu}\)、\(\beta:C_{vu}\to C_v\)，满足
\[
 \alpha i_u=i_{vu}v,\quad p_{vu}\alpha=p_u,\quad
 \beta i_{vu}=i_v,\quad p_v\beta=u[1]p_{vu},
\]
并使
\[
 C_u\xrightarrow{\alpha}C_{vu}\xrightarrow{\beta}C_v
 \xrightarrow{i_u[1]p_v}C_u[1]
\]
为好三角。这一形式称为\term[bamian-ti-gongli]{八面体公理}[octahedral axiom]，名称来自把上述四个三角画成八面体的四个面。把六个未平移对象放在同一张平面图中，并在底端补出连接箭头的目标，就能同时追踪四个三角：
\[
\begin{tikzcd}[column sep=1.2em,row sep=2.1em]
&&X\arrow[dl,"u"']\arrow[dr,"vu"]&&\\
&Y\arrow[rr,"v"]\arrow[dl,"i_u"']&&Z\arrow[dl,"i_{vu}"']\arrow[dr,"i_v"]&\\
C_u\arrow[rr,"\alpha"]\arrow[d,"p_u"']&&C_{vu}\arrow[rr,"\beta"]\arrow[dll,"p_{vu}"]&&C_v\arrow[d,"p_v"]\arrow[dll,"\gamma"]\\
X[1]\arrow[rrrr,bend right=25,"{u[1]}"']&&C_u[1]&&Y[1]\arrow[ll,"{i_u[1]}"] .
\end{tikzcd}
\]
这里 \(\gamma=i_u[1]p_v\)。顶部的两条路径比较 \(v\circ u\)；中间两条新箭头把 \(u,vu,v\) 的三个锥连成第四个三角。底端平移对象上的箭头记录各三角的连接态射，图中的交换关系正是前面四个等式。后面的锥矩阵计算逐一验证这些路径，以及第四个三角为何仍为好三角。

\begin{lemma}\label{b4:lem:pretriangle-hom}
设局部小加性范畴 \(\mathcal T\) 带平移及满足 (TR1)--(TR3) 的好三角。对好三角 \(X\xrightarrow{u}Y\xrightarrow{v}Z\xrightarrow{w}X[1]\)，相邻箭头复合为零，且对任意对象 \(W\in\mathcal T\)，两个 Hom 序列
\[
 \operatorname{Hom}(W,X)\longrightarrow\operatorname{Hom}(W,Y)
 \longrightarrow\operatorname{Hom}(W,Z),
\]
\[
 \operatorname{Hom}(Z,W)\longrightarrow\operatorname{Hom}(Y,W)
 \longrightarrow\operatorname{Hom}(X,W)
\]
均在中间正合。三角态射中任意两个分量为同构，则第三个也为同构。
\end{lemma}
\begin{proof}
从恒等三角到给定三角应用 (TR3)，前两个分量取 \(1_X,u\)，第三个分量由零对象出发，得到 \(vu=0\)。旋转后同理得到其余零复合。

若 \(a:W\to Y\) 满足 \(va=0\)，要证明的是存在 \(b_0:W\to X\)，使 \(ub_0=a\)。(TR3) 从前两项的交换方块补出第三项的态射；因此取第二项为零的好三角，使方块的交换条件恰好是 \(va=0\)，并让补出的态射经逆平移成为所需的 \(b_0\)。具体地，对上方好三角
\(W\to0\to W[1]\xrightarrow{-1}W[1]\)
和下方旋转三角
\(Y\xrightarrow{v}Z\xrightarrow{w}X[1]\xrightarrow{-u[1]}Y[1]\)
使用 (TR3)，前两个分量取 \(a,0\)。所得第三个分量 \(b:W[1]\to X[1]\) 满足 \(u[1]b=a[1]\)，所以 \(a=ub[-1]\)。这给出第一列的反向包含。

若 \(a:Y\to W\) 满足 \(au=0\)，恒等三角的逆旋转给出好三角 \(0\to W\xrightarrow{1}W\to0\)。从给定三角到这个三角应用 (TR3)，前两个分量取 \(0:X\to0\) 和 \(a:Y\to W\)；前方块由 \(au=0\) 交换，补出的第三个分量 \(b:Z\to W\) 满足 \(bv=a\)。结合 \(vu=0\)，这就证明了反变 Hom 序列在中间正合。

把两类正合列旋转延长。若一个三角态射的前两个分量 \(a:X\to X'\)、\(b:Y\to Y'\) 同构，对每个 \(W\) 比较两个三角产生的连续五项
\[
\begin{gathered}
 \operatorname{Hom}(W,X)\longrightarrow\operatorname{Hom}(W,Y)
 \longrightarrow\operatorname{Hom}(W,Z)\longrightarrow{}\\
 \operatorname{Hom}(W,X[1])\longrightarrow\operatorname{Hom}(W,Y[1]).
\end{gathered}
\]
它们在中间三项正合，五个比较映射依次由 \(a,b,c,a[1],b[1]\) 诱导。除中央由 \(c:Z\to Z'\) 诱导的映射外，其余四个都是同构，所以\cref{b4:thm:five-lemma}给出 \(\operatorname{Hom}(W,c)\) 为同构。取 \(W=Z'\)，由满射性得到 \(s:Z'\to Z\) 满足 \(cs=1_{Z'}\)；再取 \(W=Z\)，由单射性和 \(csc=c\) 得 \(sc=1_Z\)。故 \(c\) 本身同构。其他两种情形由旋转归约。
\end{proof}

因此，同一态射的两个好三角可以用 (TR3) 补出同构；但补出的同构是否唯一，是另一回事。

\subsection{三角范畴与交换范畴的区别}
\label{b4:subsec:1602:03}

交换范畴用核和余核描述一个态射，三角范畴则用第三个对象同时保留它前后次数的信息。不能把这个对象简单称作余核：即使 \(u=0\)，它也通常是 \(Y\oplus X[1]\)，而不是 \(Y\)。

\begin{proposition}\label{b4:prop:triangle-mono-splits}
设 \(\mathcal T\) 是三角范畴。\(\mathcal T\) 中每个范畴意义的单态射均分裂，每个范畴意义的满态射也均分裂。
\end{proposition}
\begin{proof}
若 \(u:X\to Y\) 单，把它补成好三角 \(X\to Y\to Z\xrightarrow{w}X[1]\)。由 \(u[1]w=0\) 和 \(u[1]\) 单得 \(w=0\)。旋转并应用\cref{b4:lem:pretriangle-hom}的反变正合列，得到
\[
 \operatorname{Hom}(Y,X)\longrightarrow\operatorname{Hom}(X,X)
 \xrightarrow{(-w[-1])^*}\operatorname{Hom}(Z[-1],X).
\]
最后的映射为零，故 \(1_X\) 提升为 \(r:Y\to X\)，即 \(ru=1_X\)。满态射的论证对偶。
\end{proof}

在交换群范畴中，单态射 \(\mathbb Z\xrightarrow{2}\mathbb Z\) 没有左逆。三角范畴却能将每个态射补成好三角，旋转和 Hom 正合性因而对单态射施加了额外的分裂约束。后面将用同一个乘二映射说明：逐次单射的复形映射，在 \(K(\mathbf{Ab})\) 中未必是单态射。

\subsection{显式锥构造与八面体的符号}
\label{b4:subsec:1602:04}

\begin{theorem}\label{b4:thm:homotopy-triangulated}
设 \(\mathcal A\) 为局部小加性范畴。以平移 \([1]\) 和\cref{b4:def:good-cone-triangle}的好三角，\(K(\mathcal A)\) 是三角范畴；\(K^b(\mathcal A),K^+(\mathcal A),K^-(\mathcal A)\) 也各自为三角范畴。
\end{theorem}
\begin{proof}
好三角按定义是与标准锥三角同构的三角，所以对同构封闭。\(\operatorname{Cone}(1_A)\) 可缩，在 \(K(\mathcal A)\) 中是零对象；因而恒等态射的标准锥三角与 \(A\xrightarrow{1}A\to0\to A[1]\) 同构。对任意 \(u\in\operatorname{Hom}_{K(\mathcal A)}(A,B)\)，选择代表它的复形映射 \(f:A\to B\)，其标准锥三角的首箭头就是 \(u=[f]\)，从而补全了 \(u\)。这验证了 (TR1) 的三个要求。

为核对 (TR2)，写 \(C_f=B\oplus A[1]\)。包含 \(i_f:B\to C_f\) 的锥在次数 \(n\) 为
\[
 B^n\oplus A^{n+1}\oplus B^{n+1},\qquad
 \mathrm{d}(b,a,b')=(\mathrm{d}b+fa+b',-\mathrm{d}a,-\mathrm{d}b').
\]
映射 \(r(b,a,b')=a\) 和 \(s(a)=(0,a,-fa)\) 给出它与 \(A[1]\) 的同伦等价。确实 \(rs=1\)，而 \(1-sr=\mathrm{d}h+h\mathrm{d}\)，其中 \(h(b,a,b')=(0,0,b)\)。记标准锥三角中的包含为 \(j:C_f\to\operatorname{Cone}(i_f)\)，最后投影为 \(q:\operatorname{Cone}(i_f)\to B[1]\)，则 \(rj=p_f\)、\(qs=-f[1]\) 严格成立，并且在 \(K(\mathcal A)\) 中
\[
 (-f[1])r=qsr=q,
\]
因为 \(sr\simeq1\)。因此 \((1_B,1_{C_f},r)\) 的三个方块都交换，它是标准锥三角到旋转三角的同构。

锥的平移同构
\[
 \operatorname{Cone}(f[r])\xrightarrow{\sim}\operatorname{Cone}(f)[r],
 \qquad (b,a)\longmapsto(b,(-1)^ra)
\]
与包含相容，把最后投影变为 \((-1)^rp_f[r]\)。所以好三角 \(A\xrightarrow{u}B\xrightarrow{v}C\xrightarrow{w}A[1]\) 平移到 \(r=-1\) 后，箭头为 \(u[-1],v[-1],-w[-1]\) 的三角仍为好三角。将它旋转两次，所得箭头为 \(-w[-1],-u,-v\)；在第三个对象上取负恒等同构，就得到逆向旋转
\[
 C[-1]\xrightarrow{-w[-1]}A\xrightarrow{u}B\xrightarrow{v}C.
\]
由此正向和逆向旋转都保持好三角，(TR2) 的两向均成立。

为核对 (TR3)，先取两个标准锥三角，设前两个分量为 \(a:A\to A'\)、\(b:B\to B'\)。同伦交换表示可以选择 \(H\) 使
\(bf-f'a=\mathrm{d}_{B'}H+H\mathrm{d}_A\)。要与包含、投影相容，第三个分量可取上三角形式
\[
 c=\begin{pmatrix}b&T\\0&a[1]\end{pmatrix}:C_f\longrightarrow C_{f'},
 \qquad T^n:A^{n+1}\longrightarrow B'^n.
\]
将它代入复形映射条件 \(\mathrm{d}_{C_{f'}}c=c\mathrm{d}_{C_f}\)，非对角块的等式为
\[
 \mathrm{d}_{B'}T+T\mathrm{d}_A=bf-f'a.
\]
因此取 \(T^n=H^{n+1}\) 就能用所给同伦修正原方块的误差，即定义
\[
 c:C_f\longrightarrow C_{f'},\qquad c(y,x)=(by+Hx,ax).
\]
第二分量的微分相容性由 \(a\) 为复形映射得到。它与包含和投影严格相容，所以补成三角态射。任意好三角先用同构换成标准锥即可。至此前三条公理成立，因而可用\cref{b4:lem:pretriangle-hom}：对固定首箭头，不同的好三角之间存在保持前两项的同构。

现在核对 (TR4)。设复形映射 \(f:A\to B\)、\(g:B\to C\)。从 \(f\) 到 \(gf\) 的前两项比较是 \((1_A,g)\)，从 \(gf\) 到 \(g\) 的比较是 \((f,1_C)\)；两个方块都严格交换，所以在刚才的公式中取 \(H=0\)，便得到前两张锥映射。第三张先由 \(p_g\) 投影到 \(B[1]\)，再由 \(i_f[1]\) 包含到 \(C_f[1]\)。于是三张映射为
\[
 \alpha(b,a)=(gb,a),\qquad
 \beta(c,a)=(c,fa),\qquad
 \gamma(c,b)=(b,0).
\]
最后一个目标为 \(C_f[1]\)。它们都是复形映射，且满足公理中四个相容等式和 \(\gamma=i_f[1]p_g\)。复合 \(\beta\alpha\) 的零伦为 \((b,a)\mapsto(0,b)\)；要将这三个箭头识别为好三角，还须把 \(\operatorname{Cone}(\alpha)\) 与 \(C_g\) 比较。

把 \(\operatorname{Cone}(\alpha)^n\) 按顺序写作
\(C^n\oplus A^{n+1}\oplus B^{n+1}\oplus A^{n+2}\)，则
\[
 \mathrm{d}(c,a,b,a')=(\mathrm{d}c+gf a+gb,\,-\mathrm{d}a+a',\,-\mathrm{d}b-fa',\,\mathrm{d}a').
\]
其中 \(a,a'\) 两个分量仍含有 \(A\) 的重复信息。令 \(t=b+fa\)，则 \(t\) 的微分分量为
\[
 -\mathrm{d}_Bb-fa'+f(-\mathrm{d}_Aa+a')
 =-\mathrm{d}_B(b+fa)=-\mathrm{d}_Bt.
\]
因此可逆坐标变换
\[
 (c,a,b,a')\longmapsto\bigl((c,b+fa),(a,a')\bigr)
\]
将微分变为两组互不混合的微分
\[
 (c,t)\longmapsto(\mathrm{d}_Cc+gt,-\mathrm{d}_Bt),\qquad
 (a,a')\longmapsto(-\mathrm{d}_Aa+a',\mathrm{d}_Aa').
\]
这正是严格的复形同构
\[
 \operatorname{Cone}(\alpha)\cong C_g\oplus\operatorname{Cone}(1_{A[1]}).
\]
第二个直和因子由\cref{b4:prop:cone-sign-compatibility}可缩，在同伦范畴中为零。取上述分解中第一因子的投影及包含，写回原坐标就是
\[
 r(c,a,b,a')=(c,b+fa),\qquad s(c,b)=(c,0,b,0).
\]
于是 \(rs=1\)、\(sr\simeq1\)，故 \(r,s\) 是同伦逆。若 \(j:C_{gf}\to\operatorname{Cone}(\alpha)\) 是包含，\(q:\operatorname{Cone}(\alpha)\to C_f[1]\) 是最后投影，则 \(rj=\beta\)、\(qs=\gamma\) 严格成立，并且在 \(K(\mathcal A)\) 中
\[
 \gamma r=qsr=q,
\]
因为 \(sr\simeq1\)。于是 \((1_{C_f},1_{C_{gf}},r)\) 的三个方块都交换，标准锥三角经此同构成为
\(C_f\xrightarrow{\alpha}C_{gf}\xrightarrow{\beta}C_g\xrightarrow{\gamma}C_f[1]\)，正是要求的好三角。给定的三个三角若非标准形式，利用前述保持首两项的同构运送所有箭头，四个相容等式随之保留。这完成八面体公理的验证。

所有构造只使用有限双积、平移和锥。它们保持两端有界、有上界或有下界，因此三个有界方向的满子范畴同样满足全部公理。
\end{proof}

由\cref{b4:thm:homotopy-triangulated}，\(K(\mathbf{Ab})\) 是三角范畴，却不是交换范畴。若它是交换范畴，把任意态射 \(f:X\to Y\) 分解为满态射 \(e:X\to I\) 和单态射 \(m:I\to Y\)，则由\cref{b4:prop:triangle-mono-splits}可取 \(s:I\to X\)、\(r:Y\to I\)，满足 \(es=1_I\)、\(rm=1_I\)。于是 \(t=sr:Y\to X\) 满足 \(ftf=f\)。对 \(\mathbb Z[0]\) 的乘二态射，这要求某个整数 \(t\) 满足 \(4t=2\)，矛盾：次数零复形的自同态同伦类就是整数，非零的链同伦无处取值。

\ProseParagraphBreak
尤其要区分“各次数均为单射”和“在同伦范畴中为单态射”。乘二 \(\mathbb Z[0]\to\mathbb Z[0]\) 逐次单，却不是同伦范畴中的单态射。令 \(C=\operatorname{Cone}(2)\)，则 \(C[-1]\to\mathbb Z[0]\) 的投影同伦类非零，而复合乘二后为零伦。这里 \(C[-1]\) 在次数零和一为 \(\mathbb Z\)，微分为 \(-2\)；投影为零伦会要求某个整数 \(t\) 满足 \(-2t=1\)，而投影复合乘二的零伦可取 \(h^1=-1\)。这是更换态射范畴后单态射性质的变化。

\begin{definition}\label{b4:def:triangle-functor}
设 \(\mathcal T,\mathcal U\) 为三角范畴。加性函子 \(F:\mathcal T\to\mathcal U\) 连同自然同构 \(\phi_X:F(X[1])\to F(X)[1]\)，若把每个好三角 \(X\xrightarrow{u}Y\xrightarrow{v}Z\xrightarrow{w}X[1]\) 送到好三角
\[
 F(X)\xrightarrow{F(u)}F(Y)\xrightarrow{F(v)}F(Z)
 \xrightarrow{\phi_XF(w)}F(X)[1],
\]
则称 \((F,\phi)\) 为\term[sanjiao-hanzi]{三角函子}[triangle functor]，也称三角范畴间的正合函子。具有三角准逆的此类函子称为三角等价。
\end{definition}

比较同构 \(\phi\) 要把 \(F(w)\) 的目标 \(F(X[1])\) 识别为 \(F(X)[1]\)，因此它也决定最后一条连接箭头及其符号。例如对平移函子 \(F=[1]\)，两者虽都等于 \(X[2]\)，好三角平移后的箭头却是 \(u[1],v[1],-w[1]\)。取 \(\phi_X=-1_{X[2]}\) 正好给出最后的 \(-w[1]\)；若取恒等比较，得到的将是 \(w[1]\)，这个三角未必仍是好三角。下章的局部化函子则采用与复形平移严格相容的比较。对于复形范畴上的逐项加性函子，有限双积及锥矩阵的保持直接给出三角函子。
